\begin{quote}\small\noindent
\textbf{About this version.} This is the second version of the report. The first (30~October 2026) covered MICA, ModelAngelo, EModelX(+AF), EMProt and E3-CryoFold. On 31~October 2026 Prof.\ Cheng asked me to add two methods, CryoAtom and CryoZeta. \emph{Both are added in this version} (Section~\ref{sec:added}) and are marked with $^\dagger$ in Table~\ref{tab:methods}. CryoAtom ran on all \NEntries{} entries. CryoZeta, whose README asks for a GPU with 32~GB, ran on the 15~GB GPUs of the free machine as far as the free GPU time reached: it has a model for \NModelCryozeta{} entries, failed on \NFailedCryozeta{} and was not run on \NNotRunCryozeta{}, so its results are for a subset of the entries that is not a random sample (Sections~\ref{sec:added} and~\ref{sec:paired}). The entries, the AF3 inputs, the scoring protocol, the measures and the statistical tests are those of the first version. The Holm correction now covers \NFamilyWord{} method pairs instead of four, so the adjusted $p$-values of the earlier methods are somewhat larger than before: for E3-CryoFold (Wilcoxon) 0.020 in the first version and \PwHolmTmEthree{} now; for ModelAngelo above 3.0~\AA{} 0.059 (Wilcoxon) and 0.041 ($t$-test) in the first version and \PwHolmTmAngeloHigh{} and \PtHolmTmAngeloHigh{} now. The $t$-test for ModelAngelo above 3.0~\AA{} no longer falls below 0.05; no other conclusion about the earlier methods changes. The new text is in Sections~\ref{sec:added} and~\ref{sec:results}; the tables and figures list CryoAtom and CryoZeta as well.
\end{quote}
